\documentclass[12pt]{article}

\usepackage[a4paper,margin=1in]{geometry}
\usepackage[T1]{fontenc}
\usepackage[utf8]{inputenc}
\usepackage{lmodern}
\usepackage{amsmath,amssymb,amsfonts}
\usepackage{booktabs}
\usepackage{longtable}
\usepackage{array}
\usepackage{tabularx}
\usepackage{enumitem}
\usepackage{hyperref}
\usepackage{setspace}
\usepackage{ragged2e}

\setstretch{1.08}
\setlength{\parskip}{0.5em}
\setlength{\parindent}{0pt}

\hypersetup{
	colorlinks=true,
	linkcolor=black,
	citecolor=black,
	urlcolor=blue,
	pdftitle={Proof Notes for the General Theory of Cognitive Structuring},
	pdfauthor={Kostiantyn Osmolovskyi}
}

\title{Proof Notes\\
	\small for the General Theory of Cognitive Structuring}
\author{
	Kostiantyn Osmolovskyi\\
	\small Independent Researcher, Ukraine\\
	\small ORCID: 0009-0006-3144-7237\\
	\small \texttt{constantinosmol@gmail.com}
}

\date{}

% ---------------------------------------------------------------------------
% Part of a series: General Theory of Cognitive Structuring (GTCS)
% Autor: Kostiantyn Osmolovskyi (ORCID: 0009-0006-3144-7237)
% License: Creative Commons Attribution 4.0 International (CC BY 4.0)
% Repository: https://github.com/k-osmolovskyi/k-osmolovskyi.github.io
% DOI: https://doi.org/10.5281/zenodo.19705809
% ---------------------------------------------------------------------------

\begin{document}
	\maketitle

\begin{center}
	\small
	Verification Report No. 2026-5\\
	Version 1.0
\end{center}

\vspace{0.5em}
	
	\section{Purpose of this document}
	
	This document provides author-side proof notes for the \emph{General Theory of Cognitive Structuring}.
	Its purpose is to record, in a compact and operationally useful form, how the main formal results
	of the series are expected to be checked, what their minimal proof logic is, what typical risks
	of misreading or false strengthening arise, and what later proof-compendium work may still be needed.
	
	This note is not a substitute for the source papers, and it is not yet a full proof compendium.
	It functions as an intermediate canonical layer between:
	\begin{itemize}[leftmargin=2em]
		\item the \emph{Collected Propositions and Theorems},
		\item the \emph{Proof Status Register},
		\item and a future full \emph{Proof Compendium}.
	\end{itemize}
	
	Its guiding questions are:
	\begin{itemize}[leftmargin=2em]
		\item What exactly should be checked for each result?
		\item What is the expected result in normalized form?
		\item What is the minimal proof idea?
		\item What is the main risk of false strengthening, collapse of levels, or misinterpretation?
		\item What remains missing for later proof-compendium treatment?
	\end{itemize}
	
	\section{Scope and limits}
	
	This document is deliberately intermediate in status.
	
	It does \emph{not} attempt to:
	\begin{itemize}[leftmargin=2em]
		\item provide full formal proofs for every result,
		\item replace theorem statements in the source papers,
		\item settle all notational variants across the series,
		\item convert architectural claims into theorem-grade results prematurely.
	\end{itemize}
	
	It aims to:
	\begin{itemize}[leftmargin=2em]
		\item make proof logic explicit for the main result families,
		\item separate theorem-ready results from sketch-level results,
		\item preserve the intended scope of each result,
		\item and prepare later canonical proof extraction.
	\end{itemize}
	
	\section{How to read these proof notes}
	
	Each proof note is organized with the following fields:
	
	\begin{itemize}[leftmargin=2em]
		\item \textbf{Result ID}
		\item \textbf{Title}
		\item \textbf{Status}
		\item \textbf{What to check}
		\item \textbf{Expected result}
		\item \textbf{Minimal proof idea}
		\item \textbf{Dependencies}
		\item \textbf{Assumptions}
		\item \textbf{Typical risk}
		\item \textbf{Missing for compendium}
		\item \textbf{Source}
	\end{itemize}
	
	\section{Status legend}
	
	\begin{longtable}{>{\RaggedRight\arraybackslash}p{0.10\textwidth}
			>{\RaggedRight\arraybackslash}p{0.75\textwidth}}
		\toprule
		\textbf{Code} & \textbf{Meaning} \\
		\midrule
		\endfirsthead
		\toprule
		\textbf{Code} & \textbf{Meaning} \\
		\midrule
		\endhead
		
		PR & Proof-ready. Source proof is already strong enough for direct canonical extraction. \\
		
		SK & Sketch-level. Source support is substantial, but proof normalization is still needed. \\
		
		CN & Canonically normalized. Result is stable in content, but the proof layer is distributed across papers or still partly implicit. \\
		
		LC & Local-condition result. The result is valid under assumptions local to a particular branch and should not be silently globalized. \\
		
		PD & Proof-development target. A priority candidate for later proof-compendium treatment. \\
		
		\bottomrule
	\end{longtable}
	
	\section{Geometric proof notes}
	
	\subsection*{G1.01. Coherence as distance to the stability region}
	\textbf{Status:} PR \\
	\textbf{What to check:} Verify that coherence is introduced as a geometric quantity over configuration space and that the defining relation is a distance-to-region relation rather than a semantic or phenomenological similarity notion. \\
	\textbf{Expected result:} 
	\[
	C_I(x_t)=\mathrm{dist}(x_t,U_X(I_t))
	\]
	or, in earlier form,
	\[
	C(x_t)=\mathrm{dist}(x_t,U).
	\]
	\textbf{Minimal proof idea:} Once configuration space and stability region are fixed, coherence is determined by the metric relation between a point and a subset. The main task is not theorem proving in a deep sense but controlled formal definition together with preservation of the distinction between configuration, region, and metric. \\
	\textbf{Dependencies:} configuration space; stability region; distance relation. \\
	\textbf{Assumptions:} admissible distance structure on the relevant space. \\
	\textbf{Typical risk:} Collapsing coherence into semantic consistency, subjective harmony, or local tension. \\
	\textbf{Missing for compendium:} Little. Mostly notation normalization from $U$ to $U_X(I_t)$. \\
	\textbf{Source:} \emph{Coherence Evaluation in Cognitive Systems}.
	
	\subsection*{G1.02. Coherence is distinct from structural tension}
	\textbf{Status:} CN \\
	\textbf{What to check:} Verify that the series does not identify geometric coherence with the regulatory quantity $T_t$, even if local comparative relations are discussed. \\
	\textbf{Expected result:} Coherence is a geometric relation to admissible structure; structural tension is a regulatory burden composed of incompatibility and coordination cost. \\
	\textbf{Minimal proof idea:} Show that the two quantities differ in formal role, domain, and architectural level. Coherence is defined over configuration-to-region relation, whereas tension is defined over regulatory burden. Therefore identity of the two objects would be a category mistake. \\
	\textbf{Dependencies:} geometric coherence; tension decomposition. \\
	\textbf{Assumptions:} none beyond controlled separation of levels. \\
	\textbf{Typical risk:} Upgrading local comparative bounds into full formal identity between coherence and tension. \\
	\textbf{Missing for compendium:} A concise normalization note clarifying any local bounding results. \\
	\textbf{Source:} \emph{Coherence Evaluation in Cognitive Systems}.
	
	\section{Regulatory and overload proof notes}
	
	\subsection*{R2.01. Structural admissibility depends on the joint regulatory state}
	\textbf{Status:} PR \\
	\textbf{What to check:} Verify that admissibility is not defined over instantaneous pressure alone, but over at least $(T,L)$ and additional control parameters where relevant. \\
	\textbf{Expected result:} Structural updating becomes admissible only as a function of the joint regulatory state, typically
	\[
	A(T,L,R,M,\bar h).
	\]
	\textbf{Minimal proof idea:} Show that the regulatory state space is explicitly constructed so that admissibility depends on both current burden and retained overload. Then demonstrate that the stability region is defined over this joint space rather than over $T$ alone. \\
	\textbf{Dependencies:} structural tension; overload memory; admissibility operator. \\
	\textbf{Assumptions:} admissibility framework of structural admissibility paper. \\
	\textbf{Typical risk:} Treating admissibility as a one-variable threshold on current tension. \\
	\textbf{Missing for compendium:} Minimal extraction only. \\
	\textbf{Source:} \emph{Structural Admissibility in Cognitive Systems}.
	
	\subsection*{R2.02. Overload memory is not reducible to instantaneous tension}
	\textbf{Status:} PR \\
	\textbf{What to check:} Verify that the overload variable $L_t$ is introduced through a retained update rule and that distinct trajectories with the same present $T_t$ can yield different $L_t$. \\
	\textbf{Expected result:} There is no function $f$ such that
	\[
	L_t=f(T_t)
	\]
	for all admissible trajectories. \\
	\textbf{Minimal proof idea:} Construct or compare trajectories with matching instantaneous tension but different histories of excess load. Because $L_t$ depends on accumulated overload and decay dynamics, its value cannot be determined by present tension alone. \\
	\textbf{Dependencies:} overload excess; retained overload update. \\
	\textbf{Assumptions:} nontrivial history dependence and admissible trajectory variation. \\
	\textbf{Typical risk:} Mistaking $L_t$ for a smoothed version of $T_t$. \\
	\textbf{Missing for compendium:} None beyond polished canonical notation. \\
	\textbf{Source:} \emph{Overload Formation in Cognitive Processing}.
	
	\subsection*{R2.03. Regulatory state is trajectory-dependent}
	\textbf{Status:} PR \\
	\textbf{What to check:} Verify that equal present tension does not guarantee equal admissibility state if overload histories differ. \\
	\textbf{Expected result:} If
	\[
	T_t^{(1)}=T_t^{(2)}
	\quad \text{but} \quad
	L_t^{(1)}\neq L_t^{(2)},
	\]
	then admissibility states may differ. \\
	\textbf{Minimal proof idea:} Since admissibility depends jointly on current tension and retained overload, any pair of trajectories with equal $T$ but unequal $L$ yields non-equivalent regulatory states. \\
	\textbf{Dependencies:} overload memory; admissibility operator. \\
	\textbf{Assumptions:} joint-state admissibility. \\
	\textbf{Typical risk:} Treating trajectory dependence as merely narrative rather than as a formal consequence of state augmentation. \\
	\textbf{Missing for compendium:} Very little. \\
	\textbf{Source:} \emph{Trajectory-Dependent Regulation in Cognitive Systems}.
	
	\subsection*{R2.04. Overload deforms the admissibility boundary and introduces hysteresis}
	\textbf{Status:} SK / LC / PD \\
	\textbf{What to check:} Verify that the critical boundary $T_{\mathrm{crit}}(L)$ depends monotonically on overload memory and that this dependence implies path-sensitive admissibility. \\
	\textbf{Expected result:} Increasing $L$ lowers admissible tolerance for current tension, so equal present $T$ may lie inside or outside the admissible region depending on path history. \\
	\textbf{Minimal proof idea:} Use monotonic deformation of the critical boundary with respect to $L$. Then show that admissibility is path-sensitive because the same instantaneous point in tension coordinates corresponds to distinct positions relative to the moving boundary when $L$ differs. \\
	\textbf{Dependencies:} critical boundary; overload memory; admissibility geometry in $(T,L)$-space. \\
	\textbf{Assumptions:} monotonic boundary deformation; local regularity of the boundary. \\
	\textbf{Typical risk:} Overstating hysteresis as globally proved without stating the deformation assumptions. \\
	\textbf{Missing for compendium:} Explicit normalized theorem statement with assumptions separated from consequences. \\
	\textbf{Source:} \emph{Trajectory-Dependent Regulation in Cognitive Systems}.
	
	\subsection*{R2.05. The regulatory phase space is minimally two-dimensional}
	\textbf{Status:} CN \\
	\textbf{What to check:} Verify that the phase-space formulation truly requires at least two variables and that this is not merely diagrammatic convenience. \\
	\textbf{Expected result:} The minimal phase space adequate for admissibility dynamics includes at least $(T,L)$. \\
	\textbf{Minimal proof idea:} Combine the non-reducibility of overload memory with the joint-state definition of admissibility. Since $L$ cannot be collapsed into $T$ and admissibility depends on both, the phase space must have at least two dimensions. \\
	\textbf{Dependencies:} R2.01; R2.02. \\
	\textbf{Assumptions:} retained overload is a genuine state component. \\
	\textbf{Typical risk:} Treating the $(T,L)$ plane as heuristic illustration instead of minimal formal state space. \\
	\textbf{Missing for compendium:} A compact collected proof from already established results. \\
	\textbf{Source:} distributed across \emph{Structural Admissibility in Cognitive Systems}, \emph{Overload Formation in Cognitive Processing}, and \emph{Trajectory-Dependent Regulation in Cognitive Systems}.
	
	\section{Invariant and compression proof notes}
	
	\subsection*{R3.01. Invariant-induced stability geometry}
	\textbf{Status:} PR \\
	\textbf{What to check:} Verify that invariant structure is the source of admissible geometry and that the stability region is not left as an ungrounded primitive once the invariants paper is in play. \\
	\textbf{Expected result:} Given invariant structure and admissibility profiles, a stability region $U_X(I_t)$ is induced. \\
	\textbf{Minimal proof idea:} Use the admissibility conditions attached to invariants to define the admissible set of configurations. Then show that this set serves as the geometric basis for coherence evaluation. \\
	\textbf{Dependencies:} invariants; admissibility profiles; thresholds. \\
	\textbf{Assumptions:} sufficient regularity of admissibility profiles. \\
	\textbf{Typical risk:} Reading the invariants paper as if it retroactively makes the earlier coherence paper logically circular. \\
	\textbf{Missing for compendium:} Minimal normalization only. \\
	\textbf{Source:} \emph{Invariants in Cognitive Architectures}.
	
	\subsection*{R3.02. Global conflict is structured by invariant interaction architecture}
	\textbf{Status:} PR \\
	\textbf{What to check:} Verify that conflict is constructed over interaction architecture and not treated as an unstructured sum of local mismatches. \\
	\textbf{Expected result:} Global conflict load depends on the invariant interaction graph and on admissible aggregation of local incompatibilities. \\
	\textbf{Minimal proof idea:} Define pairwise incompatibility over graph edges and then apply a monotone aggregation rule. The result inherits architectural structure from the graph and cannot be reduced to independent local discrepancies. \\
	\textbf{Dependencies:} invariant structure; interaction graph; incompatibility functions. \\
	\textbf{Assumptions:} edge locality; monotone aggregation. \\
	\textbf{Typical risk:} Collapsing structural conflict into ordinary semantic contradiction or local mismatch. \\
	\textbf{Missing for compendium:} None beyond collected theorem numbering. \\
	\textbf{Source:} \emph{Invariants in Cognitive Architectures}.
	
	\subsection*{R3.03. Compression is not equivalent to simplification}
	\textbf{Status:} PR \\
	\textbf{What to check:} Verify that structural reduction alone is insufficient and that compression is defined through its stabilizing or overload-reducing role. \\
	\textbf{Expected result:} A structurally smaller organization need not count as a compression unless it satisfies the stronger compression criterion. \\
	\textbf{Minimal proof idea:} Compare mere reduction in structural size with the stricter criterion involving expected overload and admissible transformation. A counterexample schema is enough to separate the two notions. \\
	\textbf{Dependencies:} admissible updating; expected overload criterion. \\
	\textbf{Assumptions:} compression framework. \\
	\textbf{Typical risk:} Using ``compression'' as a synonym for any reduction or simplification. \\
	\textbf{Missing for compendium:} None substantial. \\
	\textbf{Source:} \emph{Structural Compression in Cognitive Architectures}.
	
	\subsection*{R3.04. Compression generates a new invariant}
	\textbf{Status:} PR \\
	\textbf{What to check:} Verify that the compressed object is not merely an eliminated trace of prior structure but enters the updated architecture as a preserved organizing condition. \\
	\textbf{Expected result:} In
	\[
	I_{t+1}=(I_t\setminus S)\cup\{F(S)\},
	\]
	the compressed object $F(S)$ functions as a new invariant. \\
	\textbf{Minimal proof idea:} Show that $F(S)$ is preserved under ordinary processing in the updated architecture and continues to constrain admissible organization. That is enough for invariant status. \\
	\textbf{Dependencies:} compression operation; invariant definition. \\
	\textbf{Assumptions:} preserved role of the compressed object at the updated level. \\
	\textbf{Typical risk:} Treating compression as eliminative rather than generative. \\
	\textbf{Missing for compendium:} Minimal. \\
	\textbf{Source:} \emph{Structural Compression in Cognitive Architectures}.
	
	\subsection*{R3.05. Compression can preserve or transform admissible geometry}
	\textbf{Status:} SK / PD \\
	\textbf{What to check:} Verify that not all compressions yield level transition and that geometric equivalence must be checked after transformation. \\
	\textbf{Expected result:} Some compressions preserve the equivalence class of admissible geometry; others transform it non-equivalently. \\
	\textbf{Minimal proof idea:} Compare pre- and post-compression admissible regions under the chosen equivalence relation. If they remain equivalent, the compression is level-preserving; if not, it is level-forming. \\
	\textbf{Dependencies:} compression operation; architectural level; geometric equivalence. \\
	\textbf{Assumptions:} admissible equivalence notion for geometric comparison. \\
	\textbf{Typical risk:} Assuming every compression automatically implies level shift. \\
	\textbf{Missing for compendium:} A compact theorem-family statement separating level-preserving from level-forming compression. \\
	\textbf{Source:} \emph{Structural Compression in Cognitive Architectures}.
	
	\subsection*{R3.06. Level transition is a non-equivalent change in admissible geometry}
	\textbf{Status:} CN / PD \\
	\textbf{What to check:} Verify that level transition is not defined by size, intensity, or complexity, but by non-equivalence of admissible geometry. \\
	\textbf{Expected result:} Level transition occurs iff admissible geometry changes non-equivalently under the adopted level criterion. \\
	\textbf{Minimal proof idea:} Combine the level definition from structural admissibility with the geometric comparison logic made operational in the compression branch. \\
	\textbf{Dependencies:} architectural level; admissible geometry; compression branch. \\
	\textbf{Assumptions:} well-defined equivalence relation over admissible geometry. \\
	\textbf{Typical risk:} Replacing geometric non-equivalence with informal notions of greater complexity. \\
	\textbf{Missing for compendium:} A single normalized collected proof spanning two papers. \\
	\textbf{Source:} \emph{Structural Admissibility in Cognitive Systems}, strengthened in \emph{Structural Compression in Cognitive Architectures}.
	
	\section{Representation and accessibility proof notes}
	
	\subsection*{R4.01. Coherence representation is an internal regulatory variable rather than a reconstruction of geometric coherence}
	\textbf{Status:} PR \\
	\textbf{What to check:} Verify that the representational variable is introduced from observable or admissible signals and not from direct access to full geometric coherence. \\
	\textbf{Expected result:} Internal regulatory representation is a proxy built from signal structure, not a full reconstruction of $C_I(x_t)$. \\
	\textbf{Minimal proof idea:} Use the bounded-access premise of the representational branch. If only partial signal structure is available, the resulting representation is necessarily proxy-like rather than globally reconstructive. \\
	\textbf{Dependencies:} geometric coherence; observable structural signals; representation-construction map. \\
	\textbf{Assumptions:} partial observability of structural stability. \\
	\textbf{Typical risk:} Treating coherence representation as if it were a direct self-model of geometry. \\
	\textbf{Missing for compendium:} None substantial. \\
	\textbf{Source:} \emph{Emergence of Coherence Representation}.
	
	\subsection*{R4.02. Coherence representation preserves regulatory order without requiring metric preservation}
	\textbf{Status:} CN / PD \\
	\textbf{What to check:} Verify that ordinal adequacy is sufficient for regulation and that metric reconstruction is not required. \\
	\textbf{Expected result:} Internal representation may preserve regulatory ranking while failing to preserve exact geometric distance. \\
	\textbf{Minimal proof idea:} Show that regulatory use only requires a stable ordering sufficient for action or regime selection. Then show that order preservation is weaker than metric preservation. \\
	\textbf{Dependencies:} coherence representation; regulatory ordering; multi-system order logic. \\
	\textbf{Assumptions:} order-sensitive regulation. \\
	\textbf{Typical risk:} Assuming representational adequacy requires exact metric fidelity. \\
	\textbf{Missing for compendium:} A collected theorem-style statement integrating single-system and multi-system sources. \\
	\textbf{Source:} \emph{Emergence of Coherence Representation}; \emph{Coherence Representation in Multi-System Regulation}.
	
	\subsection*{R4.04. Pre-symbolic admissibility determines the admissible discrepancy domain}
	\textbf{Status:} SK / PD \\
	\textbf{What to check:} Verify that only discrepancies admitted by the pre-symbolic operator enter the downstream discrepancy domain. \\
	\textbf{Expected result:} 
	\[
	D_t=Z_t^{adm}
	\]
	is determined by pre-symbolic admissibility rather than by the full discrepancy field. \\
	\textbf{Minimal proof idea:} Define the signal field before filtering, define the pre-symbolic admissibility operator, then show that the downstream admissible domain is the image of that filtering operation. \\
	\textbf{Dependencies:} pre-symbolic signal domain; pre-symbolic admissibility operator. \\
	\textbf{Assumptions:} admitted discrepancies are the only discrepancies propagated downstream. \\
	\textbf{Typical risk:} Collapsing the full discrepancy space into the admissible discrepancy domain. \\
	\textbf{Missing for compendium:} Explicit normalized proof statement with operator-domain language. \\
	\textbf{Source:} \emph{Pre-Symbolic Admissibility in Cognitive Systems}.
	
	\subsection*{R4.05. Restricted accessibility of coherence}
	\textbf{Status:} PR \\
	\textbf{What to check:} Verify that geometric deviation from admissible structure need not generate a regulatorily accessible signal if the relevant discrepancies are filtered out upstream. \\
	\textbf{Expected result:} Coherence may be geometrically nonzero while no corresponding regulatory signal is generated. \\
	\textbf{Minimal proof idea:} Combine geometric coherence with the admissible discrepancy domain and show that filtering can block transmission from geometry to accessible signal. \\
	\textbf{Dependencies:} geometric coherence; admissible discrepancy domain; regulatory accessibility. \\
	\textbf{Assumptions:} pre-symbolic admissibility acts prior to stable signal formation. \\
	\textbf{Typical risk:} Assuming that every geometric incoherence must automatically be regulatorily represented. \\
	\textbf{Missing for compendium:} Minimal. \\
	\textbf{Source:} \emph{Restricted Accessibility of Coherence in Cognitive Systems}.
	
	\subsection*{R4.06. Non-injectivity of regulatory accessibility}
	\textbf{Status:} SK / PD \\
	\textbf{What to check:} Verify that distinct configurations can yield the same accessible regulatory signal structure if their differences are filtered out. \\
	\textbf{Expected result:} There exist distinct global states that are regulatorily indistinguishable. \\
	\textbf{Minimal proof idea:} Construct or characterize two configurations differing only in discrepancies excluded from the admissible discrepancy domain. Their downstream accessibility image is therefore identical or indistinguishable. \\
	\textbf{Dependencies:} restricted accessibility; admissible discrepancy domain; signal construction. \\
	\textbf{Assumptions:} nontrivial filtering of global discrepancy structure. \\
	\textbf{Typical risk:} Mistaking structural non-injectivity for representational noise or measurement error. \\
	\textbf{Missing for compendium:} A sharper canonical proposition with explicit existence conditions. \\
	\textbf{Source:} \emph{Restricted Accessibility of Coherence in Cognitive Systems}.
	
	\section{Identity and multi-system proof notes}
	
	\subsection*{R5.01. Identity is derived as a regulatory attractor rather than introduced as a primitive self-representational object}
	\textbf{Status:} CN / PD \\
	\textbf{What to check:} Verify that the identity branch derives persistence from bounded regulatory dynamics instead of positing a primitive symbolic self-core. \\
	\textbf{Expected result:} Identity is architecturally downstream of overload-sensitive regulatory dynamics. \\
	\textbf{Minimal proof idea:} Show that repeated return toward low-overload admissible regions provides the persistence structure attributed to identity, without requiring primitive self-representational ontology. \\
	\textbf{Dependencies:} overload memory; bounded regulation; attractor interpretation. \\
	\textbf{Assumptions:} long-run regulatory setting local to the identity branch. \\
	\textbf{Typical risk:} Re-importing narrative or semantic self-description as the primitive basis of identity. \\
	\textbf{Missing for compendium:} A compact theorem-family separation between architectural identity claim and long-run concentration results. \\
	\textbf{Source:} \emph{Identity as a Regulatory Attractor}.
	
	\subsection*{R5.02. Low-overload concentration regions may function as identity-cores}
	\textbf{Status:} SK / LC / PD \\
	\textbf{What to check:} Verify that concentration near low-overload regions is stated under explicit local long-run assumptions and not silently globalized. \\
	\textbf{Expected result:} Under suitable boundedness or concentration assumptions, low-overload recurrent regions may function as identity-cores. \\
	\textbf{Minimal proof idea:} Use long-run dynamics or concentration arguments over the regulatory phase space to show recurrent or attractor-like occupation of low-overload regions. \\
	\textbf{Dependencies:} overload dynamics; admissibility geometry; attractor interpretation. \\
	\textbf{Assumptions:} local long-run assumptions from the identity branch. \\
	\textbf{Typical risk:} Treating the result as unconditional across all architectures. \\
	\textbf{Missing for compendium:} Explicit assumption-indexed theorem normalization. \\
	\textbf{Source:} \emph{Identity as a Regulatory Attractor}.
	
	\subsection*{R5.05. Minimal transferable structure between systems is ordinal rather than metric}
	\textbf{Status:} PR \\
	\textbf{What to check:} Verify that the multi-system branch requires order-preserving alignment rather than shared metric geometry. \\
	\textbf{Expected result:} Inter-system comparability is grounded in ordinal transfer conditions such as $\phi_{AB}$, not in shared metric reconstruction. \\
	\textbf{Minimal proof idea:} Show that regulatory alignment only needs preservation of ordering relevant to action or coordination. Metric identity is stronger than necessary. \\
	\textbf{Dependencies:} coherence representation; order-preserving map. \\
	\textbf{Assumptions:} ordinal adequacy of the representational layer. \\
	\textbf{Typical risk:} Assuming shared geometry is required for inter-system comparability. \\
	\textbf{Missing for compendium:} Minimal. \\
	\textbf{Source:} \emph{Coherence Representation in Multi-System Regulation}.
	
	\subsection*{R5.07. Inter-system conflict is incompatibility of admissible discrepancy domains sufficient for coordinated regulation}
	\textbf{Status:} SK / PD \\
	\textbf{What to check:} Verify that conflict is defined at the level of admissible discrepancy structure and not reduced to representational disagreement. \\
	\textbf{Expected result:} Conflict obtains when no shared admissible discrepancy structure sufficient for coordination is available. \\
	\textbf{Minimal proof idea:} Define admissible discrepancy domains for each system and then formulate compatibility as existence of sufficient overlap or transferable structure. Conflict is the failure of that condition. \\
	\textbf{Dependencies:} admissible discrepancy domains; multi-system interaction criterion. \\
	\textbf{Assumptions:} admissibility-based definition of coordination. \\
	\textbf{Typical risk:} Reducing conflict to semantic contradiction or explicit disagreement alone. \\
	\textbf{Missing for compendium:} A sharpened theorem-grade separation between definitional and derived parts of the conflict branch. \\
	\textbf{Source:} \emph{Inter-System Conflict Geometry in Cognitive Systems}.
	
	\section{Priority proof-development targets}
	
	The following results are the highest-priority candidates for later proof-compendium work:
	
	\begin{itemize}[leftmargin=2em]
		\item R2.04 --- Overload deforms the admissibility boundary and introduces hysteresis
		\item R3.05 --- Compression can preserve or transform admissible geometry
		\item R3.06 --- Level transition is a non-equivalent change in admissible geometry
		\item R4.02 --- Coherence representation preserves regulatory order without requiring metric preservation
		\item R4.04 --- Pre-symbolic admissibility determines the admissible discrepancy domain
		\item R4.06 --- Non-injectivity of regulatory accessibility
		\item R5.01 --- Identity is derived as a regulatory attractor rather than introduced as a primitive self-representational object
		\item R5.02 --- Low-overload concentration regions may function as identity-cores
		\item R5.07 --- Inter-system conflict is incompatibility of admissible discrepancy domains sufficient for coordinated regulation
	\end{itemize}
	
	\section{Compact development table}
	
	\begin{longtable}{>{\RaggedRight\arraybackslash}p{0.12\textwidth}
			>{\RaggedRight\arraybackslash}p{0.46\textwidth}
			>{\RaggedRight\arraybackslash}p{0.12\textwidth}
			>{\RaggedRight\arraybackslash}p{0.18\textwidth}}
		\toprule
		\textbf{ID} & \textbf{Result} & \textbf{Status} & \textbf{Priority} \\
		\midrule
		\endfirsthead
		\toprule
		\textbf{ID} & \textbf{Result} & \textbf{Status} & \textbf{Priority} \\
		\midrule
		\endhead
		
		G1.01 & Coherence as distance to the stability region & PR & low \\
		
		R2.01 & Structural admissibility depends on the joint regulatory state & PR & low \\
		
		R2.02 & Overload memory is not reducible to instantaneous tension & PR & low \\
		
		R2.03 & Regulatory state is trajectory-dependent & PR & low \\
		
		R2.04 & Overload deforms the admissibility boundary and introduces hysteresis & SK / PD & high \\
		
		R3.05 & Compression can preserve or transform admissible geometry & SK / PD & high \\
		
		R3.06 & Level transition is a non-equivalent change in admissible geometry & CN / PD & high \\
		
		R4.02 & Coherence representation preserves regulatory order without requiring metric preservation & CN / PD & medium-high \\
		
		R4.04 & Pre-symbolic admissibility determines the admissible discrepancy domain & SK / PD & medium-high \\
		
		R4.06 & Non-injectivity of regulatory accessibility & SK / PD & high \\
		
		R5.01 & Identity is derived as a regulatory attractor rather than introduced as a primitive self-representational object & CN / PD & medium-high \\
		
		R5.02 & Low-overload concentration regions may function as identity-cores & SK / LC / PD & high \\
		
		R5.05 & Minimal transferable structure between systems is ordinal rather than metric & PR & low \\
		
		R5.07 & Inter-system conflict is incompatibility of admissible discrepancy domains sufficient for coordinated regulation & SK / PD & high \\
		
		\bottomrule
	\end{longtable}
	
	\section{Closing note}
	
	These proof notes are intended to preserve author-side control over the emerging canonical proof
	layer of the \emph{General Theory of Cognitive Structuring}. They do not replace the papers, and
	they do not yet constitute a full proof compendium. Their role is to make explicit what should
	be checked, what the expected result is, where the main proof risks lie, and what remains to be
	done before full proof-compendium treatment.
	
\end{document}